% Generated draft: requires mathematical and source review.
\documentclass[11pt]{article}
\usepackage[margin=25mm]{geometry}
\usepackage{fontspec}
\setmainfont{Latin Modern Roman}
\usepackage{amsmath,amssymb,mathtools,unicode-math}
\setmathfont{Latin Modern Math}
\usepackage{enumitem,needspace,xurl,hyperref}
\setlength{\parindent}{0pt}
\setlength{\parskip}{5pt}
\setlength{\emergencystretch}{3em}
\newcommand{\N}{\mathbb{N}}
\newcommand{\Z}{\mathbb{Z}}
\newcommand{\Q}{\mathbb{Q}}
\newcommand{\R}{\mathbb{R}}
\newcommand{\C}{\mathbb{C}}
\newcommand{\F}{\mathbb{F}}
\newcommand{\A}{\mathbb{A}}
\newcommand{\PP}{\mathbb{P}}
\newcommand{\E}{\mathbb{E}}
\newcommand{\eps}{\varepsilon}
\newcommand{\dd}{\,\mathrm{d}}
\begin{document}
\section*{UnsolvedMath ID: 30004841}
\textbf{Source problem code:} catalogue entry
\subsection{Definitions and mathematical statement}
Let \(k\) denote the ground field, which is not specified in the supplied record, and let \(A\) be a Hopf algebra over \(k\), with counit \(\varepsilon:A\to k\). Define its augmentation ideal and augmentation residual by
\[
I=\ker\varepsilon,\qquad I^n=\operatorname{span}_k\{a_1\cdots a_n:a_1,\ldots,a_n\in I\}\quad(n\geq1),\qquad I^\infty=\bigcap_{n\geq1}I^n.
\]
The source describes \(I^\infty\) as a Hopf ideal and consequently intends \(A/I^\infty\) as a Hopf algebra quotient.

\begin{enumerate}
\item Characterize the Hopf algebras \(A\) for which \(I^\infty=\{0\}\). More generally, describe \(A/I^\infty\) in terms of \(A\). The ground field and any intended restrictions on the class of Hopf algebras are unresolved in the supplied record; the surrounding discussion concerns compact quantum groups.
\item Determine \(I^\infty\) when \(A=\mathcal O(O_N^+)\), the Hopf \(*\)-algebra of polynomial representative functions on the free orthogonal compact quantum group. Explicitly, for a positive integer \(N\), this is the universal unital complex \(*\)-algebra generated by self-adjoint elements \(u_{ij}\), \(1\leq i,j\leq N\), satisfying
\[
\sum_{r=1}^N u_{ir}u_{jr}=\delta_{ij}1,\qquad
\sum_{r=1}^N u_{ri}u_{rj}=\delta_{ij}1,
\]
with
\[
\Delta(u_{ij})=\sum_{r=1}^N u_{ir}\otimes u_{rj},\qquad
\varepsilon(u_{ij})=\delta_{ij},\qquad S(u_{ij})=u_{ji}.
\]
The intended range of \(N\) is not specified in the supplied record.
\item For the intended values of \(N\), decide whether there exist a discrete group \(\Gamma\) and a closed quantum subgroup inclusion \(\iota:\widehat\Gamma\hookrightarrow O_N^+\) such that \(\Gamma\) is residually torsion-free nilpotent and
\[
O_N^+=\langle SO_N,\iota(\widehat\Gamma)\rangle_{\mathrm{top}}.
\]
Here \(SO_N=\{g\in M_N(\R):g^{\mathsf T}g=1,\ \det g=1\}\) is included through its standard fundamental representation. The compact quantum group \(\widehat\Gamma\) has polynomial Hopf \(*\)-algebra \(\C[\Gamma]\), with \(\Delta(g)=g\otimes g\), \(\varepsilon(g)=1\), and \(g^*=g^{-1}\) for \(g\in\Gamma\). Residual torsion-free nilpotence means that for every \(g\in\Gamma\setminus\{e\}\), there is a homomorphism \(\phi_g:\Gamma\to H_g\) to a torsion-free nilpotent group with \(\phi_g(g)\ne e\). A group is torsion-free if \(h^m=e\), for an integer \(m\geq1\), implies \(h=e\); it is nilpotent if its lower central series terminates at the trivial group. Topological generation means that no proper closed quantum subgroup of \(O_N^+\) contains both indicated subgroups. The supplied record does not specify whether this existence assertion is intended for every positive integer \(N\), a restricted range, or a fixed \(N\).
\end{enumerate}
\subsection{Short English statement}
The problem asks when the powers of a Hopf algebra's augmentation ideal have zero intersection, how to describe the resulting quotient in general, and what this intersection is for the free orthogonal quantum group. It also asks whether that quantum group can be topologically generated by its classical special orthogonal subgroup together with the dual of a residually torsion-free nilpotent discrete group. The general ground field, any additional Hopf algebra hypotheses, and the intended range and quantification of \(N\) require clarification. The supplied partially solved label is dated curation, not an independently verified conclusion.
\subsection{Sources}
\begin{itemize}
\item Local UnsolvedMath registry, supplied complete record: permanent ID \(30004841\), title ``Separated Augmentation Filtrations and Quantum Group Generation,'' exact source problem code \texttt{catalogue entry}; collection: Oberwolfach Reports Open Problems. No individual proposer is supplied; the proposed year is recorded as \(2021\).
\item Supplied authoritative source citation: \emph{Quantum Groups -- Algebra, Analysis and Category Theory} (2022), Oberwolfach Reports, pp.\ 2397--2458, DOI \(10.4171/owr/2021/44\). The supplied OCR extraction includes a page marker for p.\ 2453. \url{https://doi.org/10.4171/owr/2021/44}.
\item Supplied literature source: ``Oberwolfach Report 44.'' \url{https://doi.org/10.4171/owr/2021/44}.
\item Supplied literature source: ``Connectedness and Gaussian parts.'' \url{https://arxiv.org/abs/2203.08030}.
\item Literature assessment dated exactly \texttt{2026-08-22}: the supplied curation labels the problem partially solved and states that the augmentation residual is characterized for classical compact groups and discrete group algebras, that strong connectedness provides useful generation criteria, and that the general compact-quantum-group problem and the final \(O_N^+\) existence question remain unresolved. This is supplied dated curation derived from the OWR record, not independently verified literature assessment; its status and cited claims require independent checking.
\item The registry labels the statement \texttt{corrected\_verified} and reports that the linked report was inspected and all three explicit questions were restored from an OCR-corrupted extraction. This verification claim is supplied provenance; no source inspection was performed for this draft.
\end{itemize}
\end{document}
